\documentclass[11pt]{report}
\usepackage[margin=1in]{geometry}
\usepackage{amsmath,amssymb,amsthm}
\usepackage{enumitem}
\usepackage{hyperref}
\usepackage{parskip}
\usepackage{booktabs}
\usepackage{array}
\hypersetup{colorlinks=true,linkcolor=black,urlcolor=blue,citecolor=black}

\theoremstyle{definition}
\newtheorem{definition}{Definition}[chapter]
\newtheorem{example}[definition]{Example}
\newtheorem{remark}[definition]{Remark}
\theoremstyle{plain}
\newtheorem{proposition}[definition]{Proposition}
\newtheorem{corollary}[definition]{Corollary}
\newtheorem{theorem}[definition]{Theorem}

\newcommand{\Uni}{\mathcal{U}}
\newcommand{\Ii}{\mathcal{I}}

\title{\textbf{Layers of Decision} \\ \large Feasibility, Relational Validity, Establishability, and Commitment \\ in History-Dependent Decision Systems}
\author{Flyxion \\ \small Independent Researcher}
\date{September 2026}

\begin{document}
\maketitle

\begin{abstract}
\noindent A decision is ordinarily modeled as a single function from a state to an action. This monograph argues that this model conflates at least four logically distinct questions --- what can be executed, what satisfies the governing relation, what the actor can establish from its own information, and what the actor should actually commit to --- into one object, and that the conflation is not a simplification but a loss of exactly the structure a decision system needs in order to fail informatively. We develop a layered alternative in which four nested, actor- and history-relative action sets, $A_i(h) \subseteq K_i(h) \subseteq C_i(h) \subseteq F_i(h)$, replace the single decision function; we show that the boundaries between these sets correspond to qualitatively different failure mechanisms rather than differences of degree; we develop a theory of observation and diagnosis as partition-refinement problems, of distributed knowledge as an evaluation-position problem rather than an inference problem, and of historical justification as a genuinely non-monotonic process in which more recorded history does not entail more currently justified action. We close with a general theory of decision adequacy of which the observational, diagnostic, and evidential cases are all instances, and a discussion of what the framework does and does not settle.
\end{abstract}
\tableofcontents

\part{The Problem of Collapsed Decisions}

\chapter{The Single-Decision Fallacy}

\section{The conventional model}
The most familiar abstraction for a decision-making system is a single function
\[
\delta: X \to \Uni
\]
mapping a state $x \in X$ directly to an action $u \in \Uni$. This representation is attractive for reasons that have nothing to do with its adequacy: it is computationally simple, it is easy to optimize against a scalar objective, it is easy to benchmark, and it presents a satisfying appearance of decisiveness --- the system always outputs \emph{an} action. None of these virtues bears on whether $\delta$ is the right kind of object to model a decision with.

\section{Four questions hidden inside one decision}
What $\delta$ actually conflates is at least four separate questions:
\begin{enumerate}
\item Can the action be executed at all?
\item Does the action satisfy the relation that governs it?
\item Can the actor \emph{establish}, from its own available information, that the relation is satisfied?
\item Given that it can be established, should the actor actually commit to it?
\end{enumerate}
We will write these as four actor- and history-relative sets, $F_i(h)$, $C_i(h)$, $K_i(h)$, and $A_i(h)$ respectively, developed in full in Part II.

\section{Why the distinctions matter}
The four questions are not four levels of confidence in the same claim; they are four different logical statuses, and an action can occupy any combination of them independently of the others. An action can be executable but relationally invalid ($u \in F_i \setminus C_i$); it can be relationally valid but epistemically uncertifiable from the actor's position ($u \in C_i \setminus K_i$); and it can be certifiable yet simply not selected by the actor's policy ($u \in K_i \setminus A_i$). Treating all three as instances of "the system did not choose $u$" erases distinctions that matter for diagnosis, for assigning responsibility, and for deciding what kind of intervention --- more resources, better rules, better evidence, or a different policy --- would actually help.

\section{Methodological commitment}
The governing principle of this monograph is stated once here and used throughout:
\begin{quote}
\emph{Do not scalarize a distinction before determining whether the distinction is structurally relevant.}
\end{quote}
A scalar score is a legitimate tool for choosing among options that are already known to be eligible. It is not a legitimate substitute for determining eligibility in the first place.

\chapter{Actions, States, Actors, and Histories}

\section{The universe of actions}
Let $\Uni$ denote the universe of possible actions. It is useful to keep several related notions distinct even though they share notation: an action \emph{type} (a description, such as "open the valve"), an action \emph{instance} (a particular token occurrence of that type), an action \emph{sequence}, and a \emph{policy} (a rule for selecting actions). Most of what follows concerns instances, but the distinction matters whenever a governing relation depends on which occurrence, in which context, is under evaluation.

\section{States}
Let $X$ be the space of physical or system states, and let
\[
T: X \times \Uni \to X
\]
be a transition function describing how a state changes under an action.

\section{Histories}
A history is a record of what has occurred up to some point, which we write either as an explicit sequence
\[
h_t = (x_0, u_0, x_1, u_1, \ldots, x_t)
\]
or, more abstractly, as an element $h_t \in H$ of a history space. The distinction between a history and the current state it terminates in is essential: two histories can agree on the current state, $x_t(h) = x_t(h')$ with $h \neq h'$, yet later actions may be evaluated differently along the two histories, because the governing relation is permitted to depend on the path taken and not merely on where that path currently stands.

\section{Actors}
Let $i \in \Ii$ index actors. Actor identity is permitted to enter the governing relation directly, as $R(i,h,u)$, rather than merely selecting which relation to apply from an actor-independent menu. This is what prepares the distinction, developed in Chapter 4, between physical equivalence of two actions and their equivalence under a governing relation: the same state transition can be relationally valid for one actor and invalid for another, or valid for the same actor under one history and not another.

\part{The Four Layers}

\chapter{Feasibility: What Can Be Done?}

\begin{definition}[Feasibility]
For actor $i$ at history $h$, the feasible set is
\[
F_i(h) \subseteq \Uni.
\]
\end{definition}

\section{Physical and technical feasibility}
Feasibility can fail for reasons that are purely physical --- unavailable resources, unreachable states, temporal constraints --- or reasons that are merely technical: an unavailable interface, an insufficient permission at the infrastructure level, incompatible formats, or exhausted computational resources. The framework does not distinguish these two sources at the level of the formalism; both are absorbed into whether the action can, as a matter of fact, be carried out.

\section{Feasibility as a relation}
Rather than treating feasibility as an intrinsic, context-free property of an action, it is more useful to define an executability relation
\[
\operatorname{Exec}(i,h,u),
\]
and then set
\[
F_i(h) = \{u \in \Uni : \operatorname{Exec}(i,h,u)\}.
\]
This keeps feasibility, like every other layer in this framework, indexed to an actor and a history rather than treated as a free-standing property of $u$ alone.

\begin{proposition}[First non-implication]
\[
u \in F_i(h) \;\not\Rightarrow\; u \in C_i(h).
\]
\end{proposition}
Executability says nothing about whether the action satisfies whatever relation governs it; that governing relation is introduced in the next chapter.

\chapter{Relational Validity: What Satisfies the Governing Relation?}

\begin{definition}[Governing relation and validity]
Let $R_i(h,u) \in \{0,1\}$ be the governing relation for actor $i$ at history $h$, evaluated on action $u$. The relationally valid set is
\[
C_i(h) = \{u \in F_i(h) : R_i(h,u)\}, \qquad\text{so } C_i(h) \subseteq F_i(h).
\]
\end{definition}

\section{Sources of governing relations}
$R_i(h,u)$ can encode invariants, explicit rules, contracts, permissions, obligations, compatibility conditions, safety conditions, resource constraints, institutional roles, or conditions that hold only between specific actors. The formalism is deliberately agnostic about which of these is in play in a given application; what matters is that whichever relation is in force is made explicit rather than folded silently into feasibility.

\section{Actor-relative validity}
In general $R(i,h,u) \neq R(j,h,u)$: the same action, at the same history, can be valid for one actor and invalid for another, because the relation can be sensitive to role, authority, or position, not merely to the action itself.

\section{Physical equivalence versus relational equivalence}
\begin{proposition}
Let $T(h,u) = T(h,v)$ denote that two actions induce the same state transition. Then
\[
T(h,u) = T(h,v) \;\not\Rightarrow\; R_i(h,u) = R_j(h,v).
\]
\end{proposition}
Two actions that are physically indistinguishable in their effect on the world can differ in relational status, both across actors and, as Chapter 20 develops further, across the same actor evaluated under different roles. This is the formal reason the framework cannot be reduced to a theory of physical state transitions alone.

\begin{proposition}[Second non-implication]
\[
u \in C_i(h) \;\not\Rightarrow\; u \in K_i(h).
\]
\end{proposition}
An action can genuinely satisfy the governing relation while the actor has no way, from its own position, of establishing that fact. This is the subject of the next chapter.

\chapter{Establishability: What Can the Actor Establish?}

\begin{definition}[Information state and observation]
Let $I_i(h)$ denote actor $i$'s local information state at history $h$; in general $I_i(h)$ is a proper part of $h$ itself. An observation map
\[
O_i: X \to Y_i
\]
(or more generally $O_i: H \to Y_i$) determines what the actor can perceive, and induces a partition of the underlying space via
\[
x_1 \sim_i x_2 \iff O_i(x_1) = O_i(x_2).
\]
\end{definition}

\begin{definition}[Establishment and certifiability]
Given an establishment predicate $E_i(I_i(h),h,u)$ --- read as "the actor has, from its own information, sufficient grounds to establish that $R_i(h,u)$ holds" --- the certifiable set is
\[
K_i(h) = \{u \in C_i(h) : E_i(I_i(h),h,u)\}, \qquad\text{so } K_i(h) \subseteq C_i(h).
\]
\end{definition}

\section{Validity versus establishment}
$u \in C_i(h)$ is a claim about the world (or about the relation governing the world); $u \in K_i(h)$ is a claim about what a particular actor, from a particular information position, can establish about the world. These are different kinds of claims, and the entire apparatus of Part IV exists to make precise the sense in which the second can fall short of the first without any defect in the actor's reasoning --- purely because of where the actor stands.

\begin{definition}[Soundness and completeness]
The inclusion $K_i(h) \subseteq C_i(h)$ is a \emph{soundness} condition: anything the actor certifies really is valid,
\[
u \in K_i(h) \Rightarrow u \in C_i(h).
\]
The converse, $C_i(h) \subseteq K_i(h)$, is a \emph{completeness} condition, and need not hold: an establishment mechanism can be perfectly sound while still being unable to certify every action that is, in fact, valid.
\end{definition}

\chapter{Commitment: What Should Actually Be Chosen?}

\begin{definition}[Commitment]
The committed set $A_i(h) \subseteq K_i(h)$ is generated by a commitment policy
\[
\pi_i: K_i(h) \times G_i \times Q_i \to \Uni,
\]
where $G_i$ represents the actor's goals or preferences and $Q_i$ represents risk or other decision criteria, so that $A_i(h) = \pi_i(K_i(h), G_i, Q_i)$.
\end{definition}

\section{Why commitment is not automatic}
When $|K_i(h)| > 1$, certification alone does not determine which action is chosen; a policy must select among the establishable candidates according to goals, preferences, and a risk functional $\rho_i(h,u)$. A common special case is
\[
u^\ast = \arg\max_{u \in K_i(h)} V_i(h,u),
\]
but the important structural point, independent of the particular policy, is that this optimization is restricted to $K_i(h)$ --- it operates \emph{after} the eligibility layers, never in place of them.

\begin{proposition}[Third non-implication]
\[
u \in K_i(h) \;\not\Rightarrow\; u \in A_i(h).
\]
\end{proposition}
In general $A_i(h) \subsetneq K_i(h)$: certification does not exhaust itself in commitment, and a policy that always committed to everything certifiable would not be a policy in any meaningful sense.

\part{The Nested Architecture}

\chapter{The Four-Set Hierarchy}

\begin{theorem}[Nested eligibility]
\[
A_i(h) \subseteq K_i(h) \subseteq C_i(h) \subseteq F_i(h).
\]
\end{theorem}

\begin{definition}[Boundaries]
\[
G^{FC}_i(h) = F_i(h) \setminus C_i(h), \qquad
G^{CK}_i(h) = C_i(h) \setminus K_i(h), \qquad
G^{KA}_i(h) = K_i(h) \setminus A_i(h).
\]
\end{definition}

\begin{center}
\begin{tabular}{@{}ll@{}}
\toprule
Difference & Interpretation \\
\midrule
$F \setminus C$ & relational exclusion \\
$C \setminus K$ & epistemic exclusion \\
$K \setminus A$ & decision-policy exclusion \\
\bottomrule
\end{tabular}
\end{center}

\begin{example}
Let $F_i(h) = \{a,b,c,d,e\}$, $C_i(h) = \{a,b,c\}$, $K_i(h) = \{a,b\}$, and $A_i(h) = \{a\}$. Then $d,e \in G^{FC}_i(h)$ fail on relational grounds, $c \in G^{CK}_i(h)$ fails on epistemic grounds, and $b \in G^{KA}_i(h)$ fails on policy grounds. These are three different diagnoses that a single scalar score, applied to $F_i(h)$ as a whole, would flatten into one undifferentiated rejection.
\end{example}

\section{Strict inclusions and degenerate cases}
The inclusions $A \subsetneq K \subsetneq C \subsetneq F$ are generally strict, and it is instructive to examine what it means for any of them to collapse to equality: $C = F$ describes a system with no operative relational constraint beyond feasibility itself; $K = C$ describes an actor with complete establishment power, able to certify every valid action from its own information; $A = K$ describes a policy that commits to everything it can certify, forgoing discretion entirely. None of these collapses should be assumed by default, and each is a substantive claim about the system under study, to be argued for rather than adopted as a modeling convenience.

\chapter{Non-Implication Calculus}

The framework's central results are most usefully organized as a catalogue of tempting but invalid inferences.

\begin{proposition}[Feasibility]
$u \in F_i(h) \not\Rightarrow u \in C_i(h)$.
\end{proposition}
\begin{proposition}[Validity]
$u \in C_i(h) \not\Rightarrow u \in K_i(h)$.
\end{proposition}
\begin{proposition}[Establishment]
$u \in K_i(h) \not\Rightarrow u \in A_i(h)$.
\end{proposition}
\begin{proposition}[Historical accumulation]
$h_t \subseteq h_{t+1} \not\Rightarrow K_i(h_t) \subseteq K_i(h_{t+1})$.
\end{proposition}
\begin{proposition}[Physical equivalence]
$T(h,u) = T(h,v) \not\Rightarrow R_i(h,u) = R_j(h,v)$.
\end{proposition}
\begin{proposition}[Diagnostic equivalence]
$D(x_1) = D(x_2) \not\Rightarrow Q(x_1) = Q(x_2)$.
\end{proposition}

Each of these will be developed and illustrated in the chapters that follow; taken together they form the monograph's central catalogue of structural errors --- the specific, avoidable mistake of inferring a stronger conclusion than the premises support, in six recurring shapes.

\part{Information and Observation}

\chapter{Observation as a Partition}

\begin{definition}[Induced partition]
An observation map $O_i: X \to Y_i$ induces an equivalence relation
\[
x_1 \sim_{O_i} x_2 \iff O_i(x_1) = O_i(x_2).
\]
\end{definition}

\section{Information as partition refinement}
We say $O_i' \succeq O_i$ when $O_i'$ induces a strictly finer partition of $X$ --- equivalently, when the partition ordering satisfies $\mathcal{P}(O_i') \preceq \mathcal{P}(O_i)$ under the refinement order on partitions. It is essential to distinguish \emph{more observations}, in the sense of a higher-dimensional or higher-volume signal, from \emph{more relevant distinctions}, in the sense of a partition that actually separates states requiring different downstream treatment. Chapter 14 returns to this distinction under the heading of diagnostic design.

\chapter{The Two Information Gaps}

\section{Gap I: insufficient information at a fixed position}
Suppose $O_i(x_1) = O_i(x_2)$ although the two states in fact require different treatment. The first, simplest repair is to refine the observer in place: replace $O_i$ with $O_i'$ such that $O_i'(x_1) \neq O_i'(x_2)$, without changing who is doing the observing.

\section{Gap II: distributed information}
A structurally different situation arises when
\[
O_i(x_1) = O_i(x_2) \quad\text{and}\quad O_j(x_1) = O_j(x_2),
\]
so that \emph{neither} observer alone can distinguish the states, yet the composed observation
\[
O_{ij} = (O_i, O_j)
\]
satisfies $O_{ij}(x_1) \neq O_{ij}(x_2)$. Here the relevant repair is not "think harder from position $i$"; no amount of refinement internal to $O_i$ can recover a distinction that simply is not present in $i$'s information channel. The gap can only be closed by changing or composing the evaluation position itself, producing a genuinely new knowledge object $K_{ij}(h)$ distinct from $K_i(h)$.

\begin{proposition}[Impossibility under a fixed position]
If a distinction is not present in $I_i$, no deterministic procedure operating only on $I_i$ can recover it.
\end{proposition}

\chapter{Local Knowledge and Distributed Knowledge}

\begin{definition}[Collective information and establishment]
For a coalition $S \subseteq \Ii$, let $I_S(h) = (I_i(h))_{i \in S}$ denote the actors' combined information, and let $K_S(h)$ denote the set of actions establishable from that joint information.
\end{definition}

\section{Observer composition}
Under suitable conditions, $K_i(h) \subseteq K_{ij}(h) \subseteq K_{ijk}(h)$: adding observers can only help, never hurt, establishment. But this is conditional, not automatic. Evidence composition, written $e_{ij} = e_i \oplus e_j$, requires a consistency condition on how individual evidence combines, and in general
\[
K_i(h) \cup K_j(h) \;\neq\; K_{ij}(h).
\]
The union of what each observer can separately establish is not, in general, the same as what the two can establish jointly: joint establishment may require genuinely joint inference over the combined evidence, not merely a pooling of separately certified conclusions. This is the precise sense in which epistemic composition can require relational composition: it is not enough for information to be additional: it must be relevant, compatible with what is already held, and actually composable into a single inference.

\part{Diagnostics and Error}

\chapter{Diagnostics as Observation}

The same partition-theoretic apparatus developed for observation in Part IV governs technical diagnostics without modification, once the relevant objects are relabeled. Let $X$ be a fault space (the set of possible underlying causes) and let $D: X \to Y$ be a diagnostic map producing observations $Y$. As before, $D$ induces an equivalence relation $x_1 \sim_D x_2 \iff D(x_1) = D(x_2)$, and the partition $X/{\sim_D}$ collects fault classes that the diagnostic system cannot tell apart.

\chapter{The Diagnostic Identifiability Problem}

\begin{definition}
Let $R: X \to \Uni$ be the correct intervention or correction as a function of the true underlying fault.
\end{definition}

\begin{theorem}[Diagnostic impossibility]
Suppose $D(x_1) = D(x_2)$ but $R(x_1) \neq R(x_2)$. Then for every correction rule $g: Y \to \Uni$,
\[
g(D(x_1)) = g(D(x_2)),
\]
and therefore $g(D(x_1)) \neq R(x_1)$ or $g(D(x_2)) \neq R(x_2)$: no correction rule composed with $D$ can be correct on both states.
\end{theorem}

This is the diagnostic instance of a general statement: \emph{a decision rule cannot distinguish states that its own representation has already identified.} A diagnostic map $D$ is \emph{sufficient} for $R$ precisely when $D(x_1) = D(x_2) \Rightarrow R(x_1) = R(x_2)$; when sufficiency fails, the defect lies in $D$, and no amount of optimizing the downstream controller $g$ can compensate for it. This distinction between representation ($D$) and controller ($g$) recurs throughout the monograph and is one of its more actionable design consequences (Chapter 39).

\chapter{Partition Refinement}

\section{Coarse versus fine diagnostics}
Write $D': X \to Y'$ for a candidate refinement of $D: X \to Y$, and say $D' \succeq D$ when $D'$'s induced partition refines $D$'s. Not every refinement matters: the relevant condition is not that $D'$ distinguish more states in general, but that it distinguish states specifically where $R$ requires different treatment,
\[
R(x_1) \neq R(x_2) \;\Rightarrow\; D'(x_1) \neq D'(x_2).
\]

\section{Minimal sufficient observation}
Among all diagnostic maps sufficient for $R$, the most useful is typically the \emph{coarsest} one, since it discards exactly the distinctions $R$ does not need while preserving all the ones it does. This reframes diagnostic design as a problem of equivalence-class engineering: the goal is not to maximize the amount of information collected, but to engineer a partition of the fault space that exactly tracks the distinctions the downstream response requires --- no coarser, and no finer than necessary.

\part{Time and History}

\chapter{Histories as Growing Records}

Histories are typically assumed to grow monotonically as a record, $h_t \subseteq h_{t+1}$: nothing already recorded is forgotten. This is a claim about the record, and it must not be confused with a claim about the current state: two histories can agree on the current state, $x_t(h) = x_t(h')$, while disagreeing on relational status, $C_i(h) \neq C_i(h')$, precisely because the governing relation is written as $R_i(h,u)$ --- a function of the path --- rather than $R_i(x,u)$, a function of the current state alone. This path-dependence is not a complication to be minimized; it is what makes the framework capable of expressing history-sensitive obligations, precedent, and accumulated commitments at all.

\chapter{Non-Monotonic Justification}

\begin{theorem}[Non-monotonicity of justification]
\[
h_t \subseteq h_{t+1} \;\not\Rightarrow\; K_i(h_t) \subseteq K_i(h_{t+1}).
\]
\end{theorem}

The reason is structural rather than incidental: the transition from $h_t$ to $h_{t+1}$ is mediated by an action $u_t$ that the actor itself took, and that action can change the governing relation under which \emph{later} actions are evaluated, so that $R_i(h_t, u) \neq R_i(h_{t+1}, u)$ for some $u$. Three qualitatively different outcomes are all possible: justification can shrink, $K_i(h_{t+1}) \subsetneq K_i(h_t)$; it can expand, $K_i(h_t) \subsetneq K_i(h_{t+1})$; or, most generally, the two sets can simply be incomparable, with $K_i(h_t) \not\subseteq K_i(h_{t+1})$ and $K_i(h_{t+1}) \not\subseteq K_i(h_t)$ both holding at once. This separates two orders that are easy to conflate: the \emph{information order} on histories, $h_t \preceq h_{t+1}$, and the \emph{decision-set order} on what is currently justified, $K_i(h_t) \subseteq K_i(h_{t+1})$. The first does not entail the second.

\chapter{Decision Systems as Dynamical Systems}

Once commitment feeds back into history, $u_t \in A_i(h_t)$ and $h_{t+1} = \Phi(h_t, u_t, \omega_t)$ for some environmental evolution $\omega_t$, the entire apparatus of the preceding chapters becomes one step in a loop:
\[
h_t \to F_i(h_t) \to C_i(h_t) \to K_i(h_t) \to A_i(h_t) \to u_t \to h_{t+1}.
\]
The actor's own action changes the future decision problem it will face: it can alter which actions are available, which permissions are held, what evidence exists, what obligations are in force, what resources remain, and what will be observable going forward. A decision architecture that treats each decision instant as an isolated classification problem, disconnected from the fact that its own output reshapes the next instant's inputs, will systematically mis-model exactly this feedback --- and it is this feedback, not merely a lag in updating beliefs, that grounds the non-monotonicity result of Chapter 16.

\part{Evidence and Certification}

\chapter{Evidence Objects}

Rather than treating $u \in K_i(h)$ as a bare boolean fact, it is useful to introduce an explicit evidence object $e_i(h,u)$ together with an evidence relation $\mathcal{E}(e,h,u)$, and a certification predicate $\operatorname{Cert}(e,h,u)$. Evidence can carry provenance, written schematically as $e = (s,t,p,m)$ for source, time, provenance, and method; two evidence objects can be composed, $e_1 \oplus e_2$, subject to a consistency condition; and two evidence objects can conflict, $e_1 \perp e_2$, in which case the effect of the conflict on $K_i(h)$ must be specified by the system rather than left implicit. The payoff of making evidence explicit is that certification becomes a triple $(h,u,e_i)$ rather than a pair $(h,u)$: something that can be inspected, transmitted, audited, composed, or invalidated, rather than a claim the system simply asserts.

\chapter{Soundness, Completeness, and Certification Failure}

\begin{definition}[Sound certification]
$\operatorname{Cert}(e,h,u) \Rightarrow R_i(h,u)$.
\end{definition}

\begin{remark}[Incomplete certification]
$R_i(h,u) \not\Rightarrow \exists e\, \operatorname{Cert}(e,h,u)$: an action can be valid without any available evidence being sufficient to certify it.
\end{remark}

If soundness itself fails --- if the system certifies $u \in K_i(h)$ for some $u \notin C_i(h)$ --- the nesting $K_i(h) \subseteq C_i(h)$ on which the entire architecture rests is broken, and every downstream guarantee built on that nesting is void. This is why soundness of the establishment mechanism, and the trust assumptions it rests on (reliability of evidence sources, integrity of provenance, validity of whatever observation model produced the evidence), deserve to be treated as first-class, explicitly stated assumptions rather than implementation details.

\part{Equivalence, Identity, and Relation}

\chapter{Physical Identity Versus Relational Identity}

Write $x_1 \equiv_T x_2$ for physical state equivalence under a chosen transition representation, and recall from Chapter 4 that $T(h,u) = T(h,v)$ does not entail $R_i(h,u) = R_i(h,v)$. It is useful to define an actor-indexed action equivalence, $u \sim_i v$, holding exactly when the governing relation treats $u$ and $v$ identically for actor $i$; this need not agree with equivalence under $T$, and the divergence between $\sim_i$ and $\equiv_T$ is precisely the content of Proposition 4.2. State-transition semantics --- describing a system solely by what physical states its actions produce --- are therefore insufficient on their own to characterize relational status.

\chapter{Relation Is Not Just Another Constraint}

Some governing relations are unary, $P(u)$, applying to an action independent of context. Others are irreducibly relational: $R(i,u)$, indexed to an actor; or $R(i,j,h,u)$, indexed to a pair of actors, a history, and an action jointly. Authority-to-action, obligation-to-role, and observer-to-evidence conditions are all of this second kind, and they resist reduction to unary physical properties of the action itself: the same physical act can be authorized for one actor-pair and not another, purely as a function of the relation between the parties, with no corresponding difference in the act. Some conditions, in other words, exist only \emph{as} relations among multiple positions and have no unary physical correlate at all.

\part{Decision Architecture}

\chapter{Why a Single Scalar Is Insufficient}

A scalar decision model $s(h,u) \in \mathbb{R}$, optimized directly over all of $\Uni$, discards the categorical structure of $F,C,K,A$: a score cannot substitute for a predicate when the predicate carries a categorical meaning that a real number cannot represent (an action is not "60\% valid"; it either satisfies the governing relation or it does not, whether or not the actor can establish which). Scalarization is not illegitimate as such --- it is exactly the right tool for selecting among candidates already known to be eligible, as in
\[
u^\ast = \arg\max_{u \in K_i(h)} V_i(h,u).
\]
Used this way, the scalar is a \emph{selection} mechanism operating strictly inside $K_i(h)$; it becomes illegitimate only when it is asked to stand in for the eligibility judgment itself.

\chapter{Typed Decision Architectures}

The four layers can be formalized as a type hierarchy, $F \hookleftarrow C \hookleftarrow K \hookleftarrow A$, with operations
\[
\operatorname{execute}: F \to \text{Outcome}, \qquad
\operatorname{validate}: F \to C, \qquad
\operatorname{establish}: C \to K, \qquad
\operatorname{commit}: K \to A.
\]
This typed reading is conceptually close to proof-carrying code: possessing an executable action is not the same \emph{type} of object as possessing an action accompanied by the evidence needed to justify it, and a well-typed architecture should make it impossible to pass an object of the wrong type --- a bare feasible action, say --- into an operation that expects a certified one.

\chapter{Architecture of a Layered Decision Engine}

An abstract implementation stages the four layers in sequence --- a feasibility engine, a relation evaluator, an evidence/establishment layer, and a commitment policy --- feeding forward from history to action to a new history, with observation, evidence, diagnostics, and evaluation-position management entering as cross-cutting systems rather than as a fifth stage in the pipeline. Two design rules follow directly from the typed reading of the previous chapter: certain shortcuts must be structurally forbidden, specifically $F \to A$ without validity and $C \to A$ without establishment; and every commitment should be traceable back to the full tuple $(h,u,e,R,\pi)$ that produced it, so that a later audit can identify not merely \emph{that} an action was committed to, but which relation it was checked against, which evidence certified it, and which policy selected it.

\part{Failure Theory}

\chapter{Taxonomy of Decision Failures}

The framework supports a taxonomy of failure considerably finer than "the system chose wrongly":
\begin{center}
\begin{tabular}{@{}p{3.6cm}p{9cm}@{}}
\toprule
Failure type & Description \\
\midrule
I --- Execution failure & $u \notin F_i(h)$: the action was never executable. \\
II --- Relational failure & $u \in F_i(h) \setminus C_i(h)$: executable but invalid. \\
III --- Epistemic failure & $u \in C_i(h) \setminus K_i(h)$: valid but uncertifiable from the actor's position. \\
IV --- Policy rejection & $u \in K_i(h) \setminus A_i(h)$: certifiable but not selected. \\
V --- Observation aliasing & $D(x_1) = D(x_2)$ with divergent required responses. \\
VI --- Position insufficiency & $K_i(h) \subsetneq K_S(h)$: a wider evaluation position would have certified more. \\
VII --- Temporal invalidation & $u \in K_i(h_t)$ but $u \notin K_i(h_{t+1})$. \\
VIII --- Relation misidentification & The system checks against $R'$ when the operative relation is $R \neq R'$. \\
\bottomrule
\end{tabular}
\end{center}
Each type calls for a different remedy --- more resources, a corrected rule, better or differently-positioned evidence, a revised policy, a finer diagnostic partition, a wider evaluation position, updated certification, or a corrected relation --- and collapsing them into a single failure category obscures which remedy is actually indicated.

\chapter{Adversarial Counterexamples}

It is worth constructing, deliberately, a minimal counterexample to every tempting implication the framework denies: an action executable but invalid ($u \in F \setminus C$); an action valid but uncertifiable ($u \in C \setminus K$); an action certifiable but unchosen ($u \in K \setminus A$); two actions with the same transition but different relational status across actors ($T(h,u)=T(h,v)$, $R_i(h,u) \neq R_j(h,v)$); a case of more history producing fewer justified actions ($h_t \subset h_{t+1}$, $K(h_{t+1}) \subset K(h_t)$); an observation whose dimensionality increases without its discriminative power increasing (more sensors, same partition); and a case where an additional observer contributes no usable knowledge because its information, though genuine, is neither relevant nor composable with what is already held. Each counterexample is small enough to state in a single finite model, and each is a direct witness to one of the non-implications catalogued in Chapter 8.

\part{Formal Results}

\chapter{Basic Propositions}

\begin{proposition}[Nested eligibility]
Given sound establishment and a commitment policy restricted to $K_i(h)$,
\[
A_i(h) \subseteq K_i(h) \subseteq C_i(h) \subseteq F_i(h).
\]
\end{proposition}

\begin{proposition}[Observation impossibility]
If $O(x_1) = O(x_2)$ and $Q(x_1) \neq Q(x_2)$, no deterministic function $g: Y \to Z$ can recover $Q$ from $O$ correctly on both states.
\end{proposition}

\begin{proposition}[Partition sufficiency]
A representation $O$ is sufficient for $Q$ if $O(x_1) = O(x_2) \Rightarrow Q(x_1) = Q(x_2)$.
\end{proposition}

\begin{proposition}[Historical non-monotonicity]
Without an invariance assumption on $C_i$ and $E_i$ across the transition $h_t \to h_{t+1}$, $h_t \subseteq h_{t+1}$ does not imply $K_i(h_t) \subseteq K_i(h_{t+1})$.
\end{proposition}

\begin{proposition}[Position enlargement]
Under suitable evidence-composition assumptions, $I_i \preceq I_S \Rightarrow K_i(h) \subseteq K_S(h)$; the composition assumptions (relevance, compatibility, consistency of $\oplus$) must be stated explicitly, since without them the inclusion can fail.
\end{proposition}

\chapter{Conditions for Monotonicity}

The preceding chapters emphasize non-monotonicity as a live possibility, not a universal law. Monotonicity, $K_i(h_t) \subseteq K_i(h_{t+1})$, can in fact be recovered under explicit conditions: if history only ever adds information without altering the operative relation, if previously valid evidence remains valid, if establishment is monotone in information, and if the evaluation position remains fixed, then more history does yield more justified action. Distinguishing \emph{non-monotonicity as a structural necessity} of the framework from \emph{non-monotonicity as a consequence of specific, checkable assumptions failing} is important: the framework does not claim justification must decrease with time, only that nothing in the bare fact of accumulating history guarantees it will not.

\chapter{Conditions for Exact Recovery}

Two distinct completeness questions are worth separating. Epistemic completeness, $K_i(h) = C_i(h)$, requires sufficiently rich observation, sound and complete evidence, stable relations, and an adequate evaluation position. Decision completeness, $A_i(h) = K_i(h)$, requires something categorically different: a commitment policy that selects \emph{every} certifiable action, which is rarely a desirable property of a policy at all, since discretion --- choosing among certified options rather than acting on all of them --- is usually the point of having a policy in the first place. Conflating these two notions of completeness is a further instance of the collapse this monograph argues against throughout.

\part{Applications}

\chapter{Autonomous Agents}

Mapping $F,C,K,A$ onto an autonomous system separates action feasibility (what the actuators and environment permit), safety and task constraints ($R_i$), perception and evidence accumulation ($O_i$, $e_i$, $K_i$), and planning/execution ($A_i$, $\pi_i$) into distinct, separately auditable stages, rather than folding perception, safety checking, and action selection into a single end-to-end policy whose failures cannot be attributed to any one of the three.

\chapter{Distributed Systems}

In a distributed system, partial observability, multiple nodes, authorization, consensus, and evidence aggregation are all instances of the distributed-knowledge apparatus of Chapter 11: the central structural fact is that $K_i \neq K_j \neq K_{ij}$ in general, so that neither a single node's local certification nor a naive union of nodes' certifications can be assumed, by default, to equal what the system as a whole can jointly establish.

\chapter{Fault Diagnosis}

Applied to fault diagnosis directly, with $X$ the space of possible faults, $D: X \to Y$ the diagnostic map, and $R: X \to \Uni$ the correct intervention, the theory of Part V governs observability, ambiguity, and fault equivalence classes, and reframes sensor design and intervention selection as the single problem of engineering a partition of $X$ that separates exactly the faults requiring different interventions.

\chapter{Security and Authentication}

Authentication is a natural, and cautionary, instance of the framework: the existence of a credential does not automatically imply authorized action, and a system can hold evidence that is entirely valid yet insufficient for a particular authorization decision. Distinguishing identity, authorization, evidence, and commitment explicitly --- rather than treating "authenticated" as a single undifferentiated status --- is precisely the discipline this monograph argues for throughout, applied here to a domain where the consequences of collapsing the distinctions are unusually visible.

\chapter{Governance and Institutional Systems}

In institutional settings, $R(i,h,u)$ depends explicitly on role or position, and the framework separates \emph{ability} to act, \emph{authority} to act, \emph{evidence} of that authority, and actual institutional \emph{commitment} to acting. Physical execution alone --- the fact that an act was carried out --- can never by itself represent institutional validity; the other three layers are doing independent work that mere execution cannot substitute for.

\chapter{Scientific and Experimental Decision Systems}

Applied to scientific reasoning, the framework distinguishes observation, inference, established result, and the decision to act on that result as four separate steps, none of which entails the next: an observation does not entail a correct inference from it, an inference does not entail that the result is established with sufficient warrant, and an established result does not entail that acting on it is the right decision, since that additionally depends on goals and risk exactly as in Chapter 6.

\part{Design Principles}

\chapter{Preserve the Boundaries}

\begin{quote}
\emph{Do not collapse $F$, $C$, $K$, $A$.}
\end{quote}
Each layer should remain separately inspectable in any implementation: a system that can report only "action rejected," with no indication of which of the three boundaries the rejection occurred at, has already discarded information the architecture was designed to preserve.

\chapter{Design for Relevant Distinguishability}

The general design criterion for any representation $f$ serving a downstream requirement $Q: X \to \Uni$ is
\[
f(x_1) = f(x_2) \;\Rightarrow\; Q(x_1) = Q(x_2),
\]
which is to say: never collapse two states into the same representation if the downstream decision needs to treat them differently. This single criterion, stated once here, is what unifies the observation theory of Part IV and the diagnostic theory of Part V; both are instances of it.

\chapter{Change the Evaluation Position When Necessary}

If $K_i(h)$ is insufficient, the default assumption should not be that better local inference will close the gap. When the missing distinction is structurally outside actor $i$'s information channel (Chapter 10, Gap II), the correct engineering response is to expand or compose the evaluation position, $K_i(h) \to K_{i,j}(h)$, rather than to invest further in improving $i$'s own reasoning over information it does not have.

\chapter{Treat Diagnostics as Partition Design}

The right question in instrumenting a system is not "how much feedback do we have," but "which states can the feedback distinguish." Diagnostic design should be approached explicitly as the engineering of the quotient $X/{\sim_D}$ with respect to the downstream response $R$, per Chapter 14, rather than as a generic exercise in maximizing signal volume or sensor count.

\chapter{Make History Explicit}

A history $h_t$ should never be modeled merely as a container of accumulated facts; it is also, and equally importantly, a determinant of the relations under which future actions will be evaluated. Writing $C_i = C_i(h)$ and $K_i = K_i(h)$ throughout this monograph, rather than $C_i$ and $K_i$ as history-independent objects, is a notational choice that encodes this design principle directly into the formalism.

\part{A General Theory of Decision Adequacy}

\chapter{Decision Adequacy}

\begin{definition}[Adequacy]
An observation system $O$ is adequate for a governing relation $R: X \times \Uni \to \{0,1\}$ if
\[
O(x_1) = O(x_2) \;\Rightarrow\; R(x_1,u) = R(x_2,u) \quad \text{for all relevant } u.
\]
\end{definition}
This is a strictly more precise requirement than merely calling an observation system "accurate": accuracy is typically a statement about $O$'s relationship to $X$ in isolation, whereas adequacy is a statement about $O$'s relationship to the specific downstream relation $R$ it is meant to serve. A perfectly accurate observation can still be inadequate, if it is accurate about the wrong thing.

\chapter{Epistemic Adequacy}

Writing $I_i(h) \models R(h,u)$ for the formal statement that actor $i$'s information entails establishment of $R(h,u)$, under whatever semantics the application calls for, opens the establishment predicate $E_i$ to the standard apparatus of evidence, uncertainty, ambiguity, observability, distributed information, and proof obligation --- the specific formal tools appropriate to a given domain, plugged into the general architecture rather than built into it.

\chapter{Relational Adequacy}

A system can have perfect observations of the wrong predicate entirely: accurate observation does not imply correct relational evaluation,
\[
\text{accurate observation} \;\not\Rightarrow\; \text{correct relational evaluation}.
\]
This is a higher-order failure than any considered so far --- not a failure of establishment given the right relation, but a failure to have represented the right relation in the first place --- and it is the formal counterpart of Failure Type VIII (Chapter 25).

\chapter{Commitment Adequacy}

Finally, $K_i(h)$ must be distinguished from the policy $\pi_i$ that generates $A_i(h)$ from it. A commitment mechanism requires its own specification --- an objective, a domain of admissible actions, a risk tolerance, a preference ordering, a tie-breaking rule, and an account of reversibility and consequences --- and this specification is not an automatic byproduct of certification. Commitment has its own mathematical theory, developed partially in Chapter 6, rather than being reducible to "select something from $K_i(h)$."

\part{Meta-Theory}

\chapter{What Counts as Information?}

The formalism has been deliberately noncommittal about which of several competing notions of information --- raw observations, partitions, $\sigma$-algebras, knowledge sets in the modal-logic sense, belief states, evidence sets, or sufficient statistics --- underlies $I_i(h)$. Each preserves the framework's core distinction (that establishment is a function of position, not merely of the world) to varying degrees and with varying formal cost, and the choice among them should be made by the application, not prescribed by the architecture.

\chapter{Truth, Knowledge, and Justification}

The framework deliberately keeps three predicates separate: $R(h,u)$, a world condition; $E_i(I_i,h,u)$, an epistemic condition; and the eventual commitment $A_i(h,u)$, a decision condition. Classical philosophical concerns with truth, knowledge, and justification map onto these three predicates respectively, but the framework takes no position on the philosophical analysis of any one of them individually; its contribution is the claim that whatever analysis is adopted, these three must not be conflated into a single status.

\chapter{Limits of the Framework}

The framework does not determine what a governing relation \emph{ought} to be, does not determine an actor's goals, does not guarantee that observations are correct, does not resolve conflicts between competing values, does not eliminate uncertainty, and does not establish that any particular commitment is morally or institutionally correct. It describes the architecture of the distinctions among feasibility, validity, establishment, and commitment; it does not supply the substantive content of any of the four.

\part{Synthesis}

\chapter{The Layered Decision Principle}

Returning to the central nesting,
\[
A_i(h) \subseteq K_i(h) \subseteq C_i(h) \subseteq F_i(h),
\]
and combining it with the observation theory $O_i: X \to Y_i$, the diagnostic theory $D: X \to Y$, and the historical dynamics $h_{t+1} = \Phi(h_t,u_t,\omega_t)$, yields a single architecture in which every layer is indexed by both an actor and a history, every epistemic claim is relative to an explicit evaluation position, and every transition from one history to the next is understood as feeding back into the relations that will govern the next round of decisions.

\chapter{The General Architecture}

The complete model routes observation and history into an information state, which determines $F_i(h)$; $F_i(h)$ narrows to $C_i(h)$ under the governing relation; $C_i(h)$ narrows to $K_i(h)$ under establishment; $K_i(h)$ narrows to $A_i(h)$ under the commitment policy; and $A_i(h)$ produces the action $u_t$ that generates $h_{t+1}$. Distributed evaluation enters laterally as $K_i(h) \to K_S(h)$; diagnostic refinement enters through the observation stage as $D \to D'$. No stage of this pipeline is optional, and no stage can be safely skipped by assuming that satisfying an earlier stage is sufficient for a later one.

\chapter{Beyond the Single Decision}

The fundamental unit of analysis this monograph proposes is not the pair $(x,u)$ of a state and an action, but the tuple
\[
(i, h, I_i, R, e, u, A),
\]
comprising the evaluation position, the historical context, the accessible information, the governing relation, the evidence, the candidate action, and the resulting commitment. Schematically,
\[
\text{history} + \text{evaluation position} + \text{information} + \text{relation} + \text{evidence} \;\to\; \text{admissible commitment}.
\]
The monograph closes on four qualifications, each a direct consequence of the chapters above, and each worth holding in mind as a standing caution against the single-decision fallacy with which it began:
\[
\text{can do} \neq \text{satisfies} \neq \text{can establish} \neq \text{should commit};
\]
\[
\text{more history} \not\Rightarrow \text{more justified action};
\]
\[
\text{more data} \not\Rightarrow \text{more relevant distinctions};
\]
\[
\text{better local inference} \not\Rightarrow \text{access to distributed facts}.
\]

\appendix

\chapter{Formal Notation}

\begin{center}
\begin{tabular}{@{}ll@{}}
\toprule
Symbol & Meaning \\
\midrule
$\Uni$ & action universe \\
$X$ & state space \\
$H$ & history space \\
$\Ii$ & actor set \\
$F_i(h)$ & feasible actions \\
$R_i(h,u)$ & governing relation \\
$C_i(h)$ & relationally valid actions \\
$I_i(h)$ & actor information state \\
$O_i$ & observation map \\
$E_i$ & establishment predicate \\
$K_i(h)$ & certifiable actions \\
$A_i(h)$ & committed actions \\
$D$ & diagnostic map \\
$T$ & state-transition map \\
$\Phi$ & history-transition map \\
$\pi_i$ & commitment policy \\
\bottomrule
\end{tabular}
\end{center}

\chapter{Set-Theoretic Results}

This appendix collects, for reference, the derivations of $C_i(h) \subseteq F_i(h)$ from the definition of $C_i$ as a subset of $F_i$ cut out by $R_i$; of $K_i(h) \subseteq C_i(h)$ from the definition of $K_i$ as a subset of $C_i$ cut out by $E_i$; and of $A_i(h) \subseteq K_i(h)$ from the requirement that the commitment policy $\pi_i$ take its values in $K_i(h)$. Each is immediate from the corresponding definition in Part II; they are gathered here as a single reference point for the central theorem of Chapter 7.

\chapter{Partition Theory}

For a function $f: X \to Y$, the induced equivalence $x_1 \sim_f x_2 \iff f(x_1) = f(x_2)$ determines a partition of $X$ into fibers of $f$. A partition $P'$ \emph{refines} $P$ when every cell of $P'$ is contained in some cell of $P$; refinement is a partial order on partitions of $X$, with the discrete partition (every point its own cell) as maximum and the trivial partition (one cell) as minimum. A representation $f$ is \emph{sufficient} for a relation $Q$ exactly when the partition induced by $f$ refines the partition induced by $Q$; this is the formal content behind every appeal to "distinguishability" in Parts IV, V, and XIV.

\chapter{Impossibility Proofs}

\begin{proof}[Proof of the diagnostic impossibility theorem, Chapter 13]
Suppose $D(x_1) = D(x_2)$ and $R(x_1) \neq R(x_2)$, and suppose toward contradiction that some $g: Y \to \Uni$ satisfies $g \circ D = R$. Then $g(D(x_1)) = R(x_1)$ and $g(D(x_2)) = R(x_2)$. But $D(x_1) = D(x_2)$ implies $g(D(x_1)) = g(D(x_2))$, so $R(x_1) = R(x_2)$, contradicting the hypothesis. Hence no such $g$ exists. \qedhere
\end{proof}
The observation-impossibility proposition of Chapter 27 is the identical argument with $D$ relabeled $O$, $R$ relabeled $Q$, and $\Uni$ relabeled $Z$; the proof does not depend on any diagnostic-specific content and is genuinely the same theorem appearing twice under different names, which is itself among the monograph's points.

\chapter{Temporal Models}

The history-transition map $h_{t+1} = \Phi(h_t, u_t, \omega_t)$ supports several regimes depending on what is assumed about $\Phi$ and about how $R_i$ depends on $h$: full monotonicity of $K_i$ obtains when $\Phi$ never alters the relation $R_i$ restricted to previously-available actions; path dependence obtains whenever $R_i(h,u)$ genuinely depends on $h$ rather than merely on the terminal state $x_t(h)$; and irreversibility obtains when some $\omega_t$ or $u_t$ removes states from future reachability entirely, a strictly stronger condition than mere non-monotonicity of $K_i$.

\chapter{Worked Formal Examples}

This appendix builds a single running scenario in six stages, each adding exactly one new piece of structure to the one before it, so that the cumulative effect of the whole apparatus can be seen concretely rather than only asserted in the abstract. The scenario is deliberately generic --- a technician, $i$, deciding whether to perform a maintenance action on a piece of equipment --- so that the reader can substitute a domain of their own choosing without altering any of the structure.

\section*{Stage 1: Bare feasibility}

Let the universe of actions be
\[
\Uni = \{a, b, c, d, e\},
\]
read informally as: $a$ = replace part, $b$ = recalibrate, $c$ = restart system, $d$ = vent pressure, $e$ = shut down permanently.

Suppose technician $i$ has the tools and access required for $a$, $b$, $c$, and $d$, but lacks the authorization credentials to perform $e$ (a permanent shutdown requires a supervisor override). Then
\[
F_i(h) = \{a,b,c,d\}.
\]
No claim has yet been made about whether any of these four actions is a good idea; $F_i(h)$ records only what the technician is physically and administratively capable of doing at all.

\section*{Stage 2: Adding a single relational constraint}

Suppose the governing relation encodes a safety rule: an action is permitted only if the equipment's pressure reading, part of the current history $h$, is below a threshold, \emph{except} for $d$ (venting), which is specifically the action designed to be used when pressure is high. Formally,
\[
R_i(h,u) = \begin{cases} \texttt{true} & \text{if } u = d, \\ \texttt{pressure}(h) < \theta & \text{otherwise.} \end{cases}
\]
Suppose at history $h_0$, $\texttt{pressure}(h_0) \geq \theta$ (the system is currently over-pressure). Then
\[
C_i(h_0) = \{u \in F_i(h_0) : R_i(h_0,u)\} = \{d\}.
\]
Every other feasible action is excluded on relational grounds: $a,b,c \in F_i(h_0) \setminus C_i(h_0)$, a Type II (relational) failure for each, even though the technician is perfectly capable of executing them. This is the first illustration of $G^{FC}_i(h)$ from Chapter 7: three actions that are feasible but not now valid, for the same underlying reason.

\section*{Stage 3: Incomplete information}

Suppose, however, that the technician's only instrument is a coarse pressure gauge with two readings, $Y_i = \{\texttt{normal}, \texttt{elevated}\}$, and that $\texttt{elevated}$ is triggered by both a genuinely over-threshold pressure state $x_1$ and a merely high-but-still-acceptable state $x_2$:
\[
O_i(x_1) = O_i(x_2) = \texttt{elevated}, \qquad \text{pressure}(x_1) \geq \theta,\ \text{pressure}(x_2) < \theta.
\]
The technician cannot, from the gauge reading alone, tell $x_1$ from $x_2$. Suppose the establishment predicate requires the actor to be certain which regime obtains before certifying $d$ as necessary (venting when not actually needed wastes a costly consumable, so it is not certified merely on suspicion):
\[
K_i(h) = \{u \in C_i(h) : E_i(I_i(h),h,u)\}, \qquad E_i(\texttt{elevated},h,d) = \texttt{false}.
\]
Then even though $d \in C_i(h_0)$ (venting really is the valid action, since the true state is $x_1$), the technician cannot establish this from the gauge alone:
\[
K_i(h_0) = \varnothing, \qquad C_i(h_0) \setminus K_i(h_0) = \{d\}.
\]
This is a Type III (epistemic) failure, and it is worth pausing on: the action the technician actually needs to take is sitting in $C_i(h_0)$, unreachable, not because anything is wrong with the technician's reasoning, but because the observation map $O_i$ is too coarse to support the establishment predicate the system demands.

\section*{Stage 4: Distributed evidence}

Suppose a second actor, $j$, has an independent instrument --- a vibration sensor --- with its own coarse readings, $Y_j = \{\texttt{quiet}, \texttt{noisy}\}$, such that
\[
O_j(x_1) = \texttt{noisy}, \qquad O_j(x_2) = \texttt{quiet}.
\]
Neither $i$ nor $j$ alone can distinguish $x_1$ from $x_2$: $j$'s reading is not ambiguous between them, but $j$'s establishment predicate for $d$ requires \emph{both} an elevated pressure reading and a noisy vibration reading, since vibration alone is a weak signal on its own. So $K_j(h_0) = \varnothing$ as well.

But the composed observation does distinguish the states:
\[
O_{ij}(x_1) = (\texttt{elevated}, \texttt{noisy}), \qquad O_{ij}(x_2) = (\texttt{elevated}, \texttt{quiet}), \qquad O_{ij}(x_1) \neq O_{ij}(x_2).
\]
If the joint establishment predicate certifies $d$ exactly when the reading is $(\texttt{elevated}, \texttt{noisy})$, then
\[
K_{ij}(h_0) = \{d\} \;\neq\; K_i(h_0) \cup K_j(h_0) = \varnothing.
\]
This is the concrete instance of Chapter 11's warning against naive aggregation: the union of what $i$ and $j$ can separately certify is empty, but their joint certification is not. The gap was closed not by either actor reasoning harder about their own instrument, but by composing two positions whose readings were individually ambiguous and jointly discriminating.

\section*{Stage 5: Diagnostic ambiguity and refinement}

Now consider the same pair of states, $x_1$ (genuine over-pressure) and $x_2$ (false alarm), reframed as a fault-diagnosis problem with fault space $X = \{x_1, x_2, x_3\}$, where $x_3$ is a third fault, a stuck valve, that happens to also trigger an elevated gauge reading. The coarse diagnostic map
\[
D(x_1) = D(x_2) = D(x_3) = \texttt{elevated}
\]
collapses all three into one observation class, while the correct interventions differ:
\[
R(x_1) = \texttt{vent}, \qquad R(x_2) = \texttt{no action}, \qquad R(x_3) = \texttt{replace valve}.
\]
By the diagnostic impossibility theorem (Chapter 13, proved in Appendix D), no correction rule $g$ composed with $D$ can be correct on all three states. Suppose an additional, more discriminating instrument is installed --- a valve-position sensor --- yielding a refined diagnostic map $D'$ with
\[
D'(x_1) = (\texttt{elevated}, \texttt{valve-ok}), \quad D'(x_2) = (\texttt{elevated}, \texttt{valve-ok}), \quad D'(x_3) = (\texttt{elevated}, \texttt{valve-stuck}).
\]
$D'$ now correctly separates $x_3$ from $\{x_1,x_2\}$, resolving the Type V (observation-aliasing) failure between the valve fault and the other two, though it still does not separate $x_1$ from $x_2$ --- that residual ambiguity is exactly the one resolved by the distributed-evidence composition of Stage 4. This illustrates that partition refinement is typically incremental: a single new instrument need not resolve every ambiguity in the fault space at once, only the specific one it is designed to discriminate.

\section*{Stage 6: A complete two-period history-dependent example}

Finally, assemble all four layers across a single time step, $h_0 \to h_1$, to exhibit the framework's sharpest result, non-monotonicity of $K_i$ (Chapter 16).

At $h_0$, combining Stages 2--4: $F_i(h_0) = \{a,b,c,d\}$, $C_i(h_0) = \{d\}$, and (using the joint position $K_{ij}$ from Stage 4) $K_{ij}(h_0) = \{d\}$. Suppose the commitment policy, facing a single certified option, commits to it: $A_{ij}(h_0) = \{d\}$, and the technician vents the system.

The action $u_0 = d$ changes the world: pressure drops below threshold, and the equipment enters a post-vent inspection regime in which the safety rule now requires a full inspection, $c$ (restart), to be preceded by $b$ (recalibration) --- a new relational requirement that did not exist at $h_0$, because it is triggered specifically by \emph{having vented}, which is now part of the history $h_1$. Formally,
\[
R_{ij}(h_1, c) = \big(\text{recalibration logged in } h_1\big),
\]
a condition that was vacuously true at $h_0$ (no venting had occurred, so the post-vent rule did not apply) and is now false, since $b$ has not yet been performed at $h_1$. Consequently
\[
C_{ij}(h_1) = \{b\}, \qquad c \notin C_{ij}(h_1),
\]
even though $c \in F_i(h_0)$ was feasible throughout and nothing about the technician's information changed. Since certification requires validity, $c \notin K_{ij}(h_1)$ either, while at $h_0$, $c$ was at least feasible and awaiting only a pressure check that never became relevant. If a different bookkeeping convention had instead certified $c$ prospectively at $h_0$ pending the vent (a modeling choice this appendix does not adopt, precisely to keep the example clean), one would obtain directly
\[
K_{ij}(h_1) \not\supseteq (\text{whatever was certified regarding } c \text{ at } h_0),
\]
the concrete shape of Chapter 16's theorem: $h_0 \subseteq h_1$, strictly more is on record, and yet the certified set has not grown monotonically alongside it --- because the technician's own prior action changed the relation governing the next one.

\chapter{Counterexample Library}

A reference library of minimal models, one per non-implication catalogued in Chapter 8 and one per adversarial case constructed in Chapter 26, is likewise a natural companion artifact: each entry should be small enough to verify by hand, and indexed against the specific proposition it witnesses.

\end{document}
